% Sections 1.1 to 1.5, from lectures/L01/appendix/a_theory.md.

\section{Talmängder}
\label{c1:sec:talmangder}
Matematiken är uppbyggd kring olika talmängder. Varje mängd är en utökning av den föregående.

\begin{booktable}{@{}p{7.6em}p{4.6em}LL@{}}
  \thead{Mängd} & \thead{Notation} & \thead{Beskrivning} & \thead{Exempel} \\
  \midrule
  Naturliga tal & $\mathbb{N}$ & Positiva heltal och noll & $0, 1, 2, 3, \ldots$ \\
  Heltal & $\mathbb{Z}$ & Naturliga tal samt negativa tal & $\ldots, -2, -1, 0, 1, 2, \ldots$ \\
  Rationella tal & $\mathbb{Q}$ & Tal som kan skrivas som $\frac{p}{q}$, där $p$ och $q$ är
    heltal och $q \neq 0$ & $\frac{1}{2},\, -\frac{3}{4},\, 0{,}75$ \\
  Irrationella tal & -- & Tal som \emph{inte} kan skrivas som $\frac{p}{q}$ &
    $\pi,\, \sqrt{2},\, e$ \\
  Reella tal & $\mathbb{R}$ & Alla rationella och irrationella tal & Samtliga tal på tallinjen \\
\end{booktable}

I elektroteknik används alla dessa talmängder. Resistansvärden är ofta rationella
($R = 4{,}7\,\text{k}\Omega$), medan konstanter som $\pi$ förekommer i formler för
växelströmssystem.

\section{Tallinjen och absolutbelopp}
\label{c1:sec:absolutbelopp}
Alla reella tal kan placeras på en \textbf{tallinje}. Talets avstånd från origo (noll) kallas
\textbf{absolutbelopp} och betecknas $\abs{a}$:
\[
  \abs{a} =
  \begin{cases}
    a & \text{om } a \geq 0 \\
    -a & \text{om } a < 0
  \end{cases}
\]

\begin{exempel}
\[
  \abs{5} = 5, \qquad \abs{-3} = 3, \qquad \abs{0} = 0
\]
\end{exempel}

Absolutbeloppet är alltid icke-negativt: $\abs{a} \geq 0$.

I elektroteknik används absolutbelopp bland annat för att ange amplituden hos en växelspänning,
oavsett om spänningen är positiv eller negativ i stunden.

\section{Räkneordning}
\label{c1:sec:rakneordning}
När ett uttryck innehåller flera räkneoperationer gäller följande prioritetsordning:
\begin{enumerate}
  \item \textbf{Parenteser}: räkna ut den innersta först och arbeta utåt.
  \item \textbf{Potenser och rötter}: beräkna exponenter och rötter.
  \item \textbf{Multiplikation och division}: från vänster till höger.
  \item \textbf{Addition och subtraktion}: från vänster till höger.
\end{enumerate}

\begin{exempel}
\[
  3 + 2 \times 4^2 - (1 + 5) = 3 + 2 \times 16 - 6 = 3 + 32 - 6 = 29
\]
\end{exempel}

\section{Bråkräkning}
\label{c1:sec:brak}

\subsection{Förenkling}
\label{c1:sec:forenkling}
Ett bråk förenklas genom att täljare och nämnare delas med sin största gemensamma delare (SGD):
\[
  \frac{12}{18} = \frac{12 \div 6}{18 \div 6} = \frac{2}{3}
\]

\subsection{Addition och subtraktion}
\label{c1:sec:bradd}
Bråk adderas och subtraheras med \textbf{gemensam nämnare}:
\[
  \frac{a}{b} + \frac{c}{d} = \frac{ad + bc}{bd}
\]

\begin{exempel}
\[
  \frac{1}{3} + \frac{1}{4} = \frac{4}{12} + \frac{3}{12} = \frac{7}{12}
\]
\end{exempel}

\subsection{Multiplikation}
\label{c1:sec:brmult}
Täljare multipliceras med täljare, nämnare med nämnare:
\[
  \frac{a}{b} \times \frac{c}{d} = \frac{ac}{bd}
\]

\begin{exempel}
\[
  \frac{2}{3} \times \frac{3}{5} = \frac{6}{15} = \frac{2}{5}
\]
\end{exempel}

\subsection{Division}
\label{c1:sec:brdiv}
Division med ett bråk är detsamma som multiplikation med bråkets inverterade värde, dess
\emph{reciprok}:
\[
  \frac{a}{b} \div \frac{c}{d} = \frac{a}{b} \times \frac{d}{c} = \frac{ad}{bc}
\]

\begin{exempel}
\[
  \frac{3}{4} \div \frac{3}{8} = \frac{3}{4} \times \frac{8}{3} = \frac{24}{12} = 2
\]
\end{exempel}

\subsection{Parallellkopplade motstånd}
\label{c1:sec:parallell}
I elektroteknik uppstår bråkräkning naturligt vid beräkning av parallellkopplade motstånd. Den
totala resistansen $R_{\text{TOT}}$ för två parallellkopplade motstånd $R_1$ och $R_2$ ges av
\[
  \frac{1}{R_{\text{TOT}}} = \frac{1}{R_1} + \frac{1}{R_2}.
\]

\begin{exempel}
Med $R_1 = 6\,\Omega$ och $R_2 = 3\,\Omega$ blir
\[
  \frac{1}{R_{\text{TOT}}} = \frac{1}{6} + \frac{1}{3} = \frac{1}{6} + \frac{2}{6} = \frac{3}{6}
  = \frac{1}{2},
\]
alltså $R_{\text{TOT}} = 2\,\Omega$.
\end{exempel}

\subsection{Serie- och parallellkoppling}
\label{c1:sec:serieparallell}
När flera motstånd sitter i samma krets kan de ersättas av ett enda motstånd, en så kallad
\textbf{ersättningsresistans}.

\textbf{Seriekoppling} innebär att motstånden är kopplade efter varandra. Ersättningsresistansen
$R_{\text{s}}$ är summan av resistanserna:
\[
  R_{\text{s}} = R_1 + R_2
\]

\textbf{Parallellkoppling} innebär att motstånden är kopplade bredvid varandra, vilket skrivs
$R_1 \parallell R_2$. Utöver formeln i \secref{c1:sec:parallell} kan parallellresistansen
$R_{\text{p}}$ av \emph{två} motstånd beräknas som produkt genom summa, vilket ofta går snabbare:
\[
  R_{\text{p}} = R_1 \parallell R_2 = \frac{R_1 \times R_2}{R_1 + R_2}
\]
Parallellresistansen blir alltid mindre än det minsta av de ingående motstånden.

\textbf{Ohms lag} knyter ihop spänningen $U$, resistansen $R$ och strömmen $I$:
\[
  U = R \times I
  \qquad \Leftrightarrow \qquad
  I = \frac{U}{R}
  \qquad \Leftrightarrow \qquad
  R = \frac{U}{I}
\]

\begin{aside}
\note[Tips] Räknas resistansen i k$\Omega$ och spänningen i V fås strömmen direkt i mA. Då slipper
man hålla reda på tiopotenser.
\end{aside}

\begin{exempel}[Förenkla en krets]
Kretsen nedan förenklas stegvis tills endast ett motstånd återstår. Först ersätts
parallellkopplingen av $R_2$ och $R_3$ med $R_{\text{p}}$, därefter adderas $R_1$ till den totala
resistansen $R_{\text{TOT}}$.

\bild[\linewidth]{lectures/L01/appendix/images/circuit_simplification.png}

Med $R_1 = 2\,\text{k}\Omega$, $R_2 = 12\,\text{k}\Omega$, $R_3 = 6\,\text{k}\Omega$ och
$U = 12\,\text{V}$ blir
\begin{gather*}
  R_{\text{p}} = \frac{12 \times 6}{12 + 6} = \frac{72}{18} = 4\,\text{k}\Omega, \\
  R_{\text{TOT}} = R_1 + R_{\text{p}} = 2 + 4 = 6\,\text{k}\Omega, \\
  I = \frac{U}{R_{\text{TOT}}} = \frac{12}{6} = 2\,\text{mA}.
\end{gather*}
Ordningen är avgörande: parallellkopplingen är en egen delberäkning som måste vara klar innan
$R_1$ adderas, precis som parenteser beräknas först enligt räkneordningen.
\end{exempel}

\section{Procent och proportion}
\label{c1:sec:procent}
\textbf{Procent} är ett annat sätt att uttrycka ett bråk med nämnaren 100:
\[
  p\,\% = \frac{p}{100}
\]
En procentandel av ett värde $A$ beräknas som
\[
  x\,\% \text{ av } A = \frac{x}{100} \times A.
\]

\begin{exempel}
Hur mycket är $30\,\%$ av $200\,\text{V}$?
\[
  \frac{30}{100} \times 200 = 60\,\text{V}
\]
\end{exempel}

\textbf{Verkningsgrad} (effektivitet) är ett vanligt procentmått i elektroteknik:
\[
  \eta = \frac{P_{\text{ut}}}{P_{\text{in}}} \times 100\,\%
\]

\begin{exempel}
En transformator tar in $500\,\text{W}$ och levererar $470\,\text{W}$. Verkningsgraden är
\[
  \eta = \frac{470}{500} \times 100\,\% = 94\,\%.
\]
\end{exempel}
